\begin{abstract}
Simultaneous speech translation (SimulST) requires incrementally producing target text as source speech arrives. Beyond unavoidable sources of latency such as ASR and decoding, two factors directly shape the read/write decision: structural word-order divergence between languages, and semantic dependency on future context not yet heard. Existing trajectory-synthesis methods commit on the basis of the observed prefix or of literal prefix agreement, but safe commitment is a counterfactual question: would this target prefix remain valid under plausible futures of the source?
We propose [METHOD], a training-data synthesis framework that operationalizes this view as distributional consensus over future-conditioned next-token distributions. Given a partial source utterance, our method samples plausible future source continuations, queries an LLM for next-token distributions conditioned on each hypothesized full source, and commits only tokens that lie in the high-probability intersection across all futures. The resulting consensus trajectories distill this expensive future-conditioned reasoning into a SimulST policy whose inference-time cost is unchanged.
We evaluate on GigaSpeech and RealSI across English$\rightarrow$Chinese, English$\rightarrow$German, and English$\rightarrow$Japanese against strong baselines including EAST, Simul-MuST-C, LA-$N$, wait-$k$, InfiniSST, and TAF.
\end{abstract}